% Generated draft: requires mathematical and source review.
\documentclass[11pt]{article}
\usepackage[margin=25mm]{geometry}
\usepackage{fontspec}
\setmainfont{Latin Modern Roman}
\usepackage{amsmath,amssymb,mathtools,unicode-math}
\setmathfont{Latin Modern Math}
\usepackage{enumitem,needspace,xurl,hyperref}
\setlength{\parindent}{0pt}
\setlength{\parskip}{5pt}
\setlength{\emergencystretch}{3em}
\newcommand{\N}{\mathbb{N}}
\newcommand{\Z}{\mathbb{Z}}
\newcommand{\Q}{\mathbb{Q}}
\newcommand{\R}{\mathbb{R}}
\newcommand{\C}{\mathbb{C}}
\newcommand{\F}{\mathbb{F}}
\newcommand{\A}{\mathbb{A}}
\newcommand{\PP}{\mathbb{P}}
\newcommand{\E}{\mathbb{E}}
\newcommand{\eps}{\varepsilon}
\newcommand{\dd}{\,\mathrm{d}}
\begin{document}
\section*{UnsolvedMath ID: 2307029}
\textbf{Source problem code:} catalogue entry
\subsection{Definitions and mathematical statement}
Let $\mathbb D=\{z\in\mathbb C:|z|<1\}$, fix $\alpha\in(0,1]$, and set $R_z=\alpha(1-|z|)$. The supplied question assumes a continuous function $f$ and the condition
\[
\int_{|\zeta-z|<R_z}f(\zeta)\,d\zeta=0
\qquad(z\in\mathbb D),
\]
and asks whether $f$ must be holomorphic on $\mathbb D$.

\textbf{Problem (L. Zalcman), pending source conventions.} Determine whether this one-radius integral condition forces analyticity, and how the answer changes when continuity is assumed only inside the disk.

The displayed notation does not define an integral: the region is an open disk but the differential is $d\zeta$, normally used for a contour integral. A Morera-type reading would replace the region by the positively oriented circle $|\zeta-z|=R_z$ and retain $d\zeta$; an area reading would retain the disk and use area measure $dA(\zeta)$. These are different conditions and neither is silently selected.

The source's first continuity assumption and its final only-in-$\mathbb D$ question are identical as printed. They may have lost a closed-disk distinction, but that correction is not supplied. This also matters when $\alpha=1$: the candidate circle can touch $\partial\mathbb D$, and at $z=0$ it is the full unit circle, where a function continuous only on the open disk is not defined. Thus the integral type and boundary hypothesis must be recovered before a complete implication can be stated. The supplied summary reports the problem open as of the 2018 edition, with no definitive resolution located.
\subsection{Short English statement}
Does a vanishing integral at one prescribed radius around each point force holomorphy? The source mixes disk and contour notation and loses the distinction between its two continuity assumptions.
\subsection{Sources}
\begin{itemize}
\item Hayman, \emph{Research Problems in Function Theory} (2018), Problem 7.29, attributed to L. Zalcman, \url{https://arxiv.org/abs/1809.07200}; source-TeX extraction, accessed 2026-07-29. The 2018 update was retained for status triage.
\item Local registry \texttt{data/unsolvedmath/problems.json}, permanent ID $2307029$, source code catalogue entry, \emph{Research Problems in Function Theory --- Problem 7.29}. The record labels this open and OPEN-TRIAGE, with NEEDS\_REVIEW source-status and rights qualifications. No separate integral-condition theorem citation is supplied.
\end{itemize}
\end{document}
